\documentclass[11pt]{article}

\usepackage[a4paper,margin=2.4cm]{geometry}
\usepackage{amsmath,amssymb}
\usepackage{booktabs}
\usepackage{siunitx}
\usepackage[hidelinks]{hyperref}
\usepackage{enumitem}

\newcommand{\Rd}{R_{\mathrm d}}
\newcommand{\Rp}{R_{\mathrm p}}

\title{Spread-Out Bragg Peak (SOBP) field design\\
for the proton-transport run}
\author{Method note}
\date{\today}

\begin{document}
\maketitle

\begin{abstract}
This note documents the method, as implemented in
\texttt{field\_design/sobp.py}, that builds the Spread-Out Bragg Peak (SOBP) of
the standard proton field. A single mono-energetic pencil deposits a sharp Bragg
peak and paints a thin track, giving an unrealistically small integral activity.
A clinical field instead spreads dose uniformly over a target volume, which is
what fixes the realistic absolute $\beta^{+}$ yield (and lets us cross-check
against Parodi \emph{et al.}\ 2008, Table~2). We construct the depth component of
such a field analytically: a weighted superposition of pristine Bragg peaks whose
weights flatten the dose over the target's depth extent, with a small attenuation
correction for the proton loss on the way in. The energies and weights are
computed in Python and written to a layer table that the Geant4 gun samples from;
the resulting plateau is verified directly against the simulated depth--dose
curve. Lateral spreading and dose normalisation are separate, implemented steps,
described at the end.
\end{abstract}

\section{Goal and scope}

We want a proton field that deposits a \textbf{flat dose plateau} across the
target's depth extent $[\Rp,\Rd]$ (proximal to distal range), so that a
prescribed dose can be painted \emph{uniformly} over the target. This note
covers the \textbf{depth} component (the SOBP proper); the lateral spreading and
the dose normalisation that complete the field are summarised in
Sec.~\ref{sec:rest}. The target geometry, its size and the prescribed dose are
all \textbf{parameters} (Sec.~\ref{sec:params}); nothing here is hard-wired to
the Parodi reference numbers.

\section{Pristine Bragg peak and the range--energy relation}
\label{sec:pristine}

A mono-energetic proton of energy $E$ stops at a well-defined depth, its
\textbf{range} $R$. To good approximation the two are related by the
Bortfeld power law~\cite{bortfeld1996},
\begin{equation}
R(E) = \alpha\,E^{\,p},
\label{eq:range}
\end{equation}
with, in water, $\alpha=\SI{2.2e-3}{cm\,MeV^{-p}}$ and $p=1.77$. For another
homogeneous material we scale the range by the (mass) stopping-power ratio,
which to first order is the inverse mass-density ratio,
$R_{\mathrm{mat}}\simeq R_{\mathrm{water}}/\rho_{\mathrm{rel}}$
($\rho_{\mathrm{rel}}\approx1.04$ for brain). The energy needed to reach a depth
$R$ is the inverse of Eq.~\eqref{eq:range}, $E(R)=(R/\alpha)^{1/p}$.

Equation~\eqref{eq:range} is used only to choose the \emph{layer energies}; the
actual peak shapes come from Geant4, and the plateau is verified there
(Sec.~\ref{sec:verify}), so small inaccuracies in $\alpha,p$ are corrected by
the verification step.

A useful consequence of Eq.~\eqref{eq:range}: a proton that has travelled to
depth $z<R$ has residual range $R-z$ and residual energy
$\propto (R-z)^{1/p}$, so its mass stopping power (dose per unit length) scales as
\begin{equation}
S(z;R)\ \propto\ \frac{\mathrm dE}{\mathrm dz}\ \propto\ (R-z)^{\,1/p-1}.
\label{eq:stop}
\end{equation}
This is the only feature of the peak shape the weight derivation needs.

\section{SOBP as a weighted superposition}

We place $N$ pristine peaks at ranges $R_i$ spanning the target depth,
\begin{equation}
R_i = \Rp + \frac{i}{N-1}\,(\Rd-\Rp),\qquad i=0,\dots,N-1,
\end{equation}
($R_0=\Rp$ proximal, $R_{N-1}=\Rd$ distal), with relative fluence weights
$w_i\ge0$. The total depth--dose is their superposition,
\begin{equation}
D_{\mathrm{SOBP}}(z)=\sum_{i\,:\,R_i\ge z} w_i\,S(z;R_i).
\label{eq:super}
\end{equation}
The gun realises Eq.~\eqref{eq:super} stochastically: each primary proton is
assigned to layer $i$ with probability $\propto w_i$ and given energy
$E_i=E(R_i)$.

\section{Analytic flattening weights}
\label{sec:weights}

We want $D_{\mathrm{SOBP}}(z)=\text{const}$ for $z\in[\Rp,\Rd]$. Treating the
range as continuous with fluence density $\Phi(R)$, Eq.~\eqref{eq:super} with
\eqref{eq:stop} becomes the Abel-type integral
\begin{equation}
D(z)=\int_z^{\Rd}\Phi(R)\,(R-z)^{1/p-1}\,\mathrm dR \stackrel{!}{=}\text{const}.
\end{equation}
This is an Abel (right-sided fractional) integral equation of order $a=1/p$;
inverting it for a constant left-hand side gives
\begin{equation}
\boxed{\;\Phi(R)\ \propto\ (\Rd-R)^{-1/p}\;}\qquad R\in[\Rp,\Rd],
\label{eq:phi}
\end{equation}
plus a finite reinforcement at the proximal edge $R=\Rp$ that fills the proximal
corner of the plateau. Since $-1/p\approx-0.565<0$, the weight rises towards the
\emph{distal} peak (which no deeper peak reinforces) and diverges integrably at
$R=\Rd$. Integrating Eq.~\eqref{eq:phi} over each range bin
$[R_i-\tfrac{\Delta}{2},\,R_i+\tfrac{\Delta}{2}]$ (with $\Delta=(\Rd-\Rp)/(N-1)$)
gives the finite discrete weights
\begin{equation}
\tilde w_i\ \propto\ (\Rd-R_i+\tfrac{\Delta}{2})^{1-1/p}-(\Rd-R_i-\tfrac{\Delta}{2})^{1-1/p},
\label{eq:wi}
\end{equation}
with both arguments clipped at zero so the distal bin stays finite (no peaks lie
beyond $\Rd$).

\subsection{Attenuation correction}
\label{sec:atten}

Equation~\eqref{eq:wi} flattens the \emph{ideal} superposition, but the
\emph{delivered} field droops from proximal to distal: protons are lost on the
way in (mainly to nuclear reactions), so deeper layers arrive depleted. We
compensate by boosting each layer by an exponential in its range,
\begin{equation}
w_i\ \propto\ \tilde w_i\,e^{\mu R_i},
\label{eq:atten}
\end{equation}
with $\mu$ [\si{cm^{-1}}] a single effective coefficient tuned so the
\emph{simulated} plateau is flat (the standard value is
$\mu=\SI{0.025}{cm^{-1}}$). The weights are then normalised, $\sum_i w_i=1$, and
stored as a probability table. The matching attenuation factor $e^{-\mu z}$
appears in the analytic depth--dose used for the sanity plot.

\paragraph{Why analytic (vs.\ fit-to-simulation).} Equations~\eqref{eq:wi} and
\eqref{eq:atten} need no simulation and are the standard, well-validated
construction. The single tuned $\mu$ absorbs the first-order delivered droop; the
verified plateau (Sec.~\ref{sec:verify}) confirms this is sufficient. The
fallback, if a flatter plateau were needed, is to replace the analytic $S(z;R_i)$
with Geant4-simulated pristine peaks and solve Eq.~\eqref{eq:super} for $w_i$ by
non-negative least squares.

\section{Parameters}
\label{sec:params}

Everything that defines the target and the field is a parameter; the standard
(Parodi reference) values are the \texttt{field\_design/sobp.py} defaults:
\begin{center}
\begin{tabular}{lll}
\toprule
symbol / flag & meaning & standard value \\
\midrule
$d_{\mathrm{prox}},d_{\mathrm{dist}}$ (\texttt{-{}-d-prox/-{}-d-dist}) & target proximal/distal depth & $5.5,\ \SI{10.5}{cm}$ \\
$\varnothing_{\mathrm t}$ & target (tumour) diameter & \SI{6}{cm} \\
$L_{\mathrm t}=d_{\mathrm{dist}}-d_{\mathrm{prox}}$ & target length (modulation width) & \SI{5}{cm} \\
$D$ & prescribed dose to the target & \SI{1}{Gy} \\
$N$ (\texttt{-{}-n-layers}) & number of energy layers & $20$ \\
$\rho_{\mathrm{rel}}$ (\texttt{-{}-rho-rel}) & material density rel.\ to water & $1.04$ (brain) \\
$\mu$ (\texttt{-{}-mu}) & attenuation-correction coefficient & \SI{0.025}{cm^{-1}} \\
material & phantom material & \texttt{G4\_BRAIN\_ICRP} \\
\bottomrule
\end{tabular}
\end{center}
The depths map to ranges $\Rp,\Rd$ in the chosen material via Sec.~\ref{sec:pristine}.

\section{Algorithm and output}

\begin{enumerate}[nosep]
\item From $(d_{\mathrm{prox}},d_{\mathrm{dist}},N,\rho_{\mathrm{rel}})$ compute
$\Rp,\Rd$ and the layer ranges $R_i$.
\item Energies $E_i=E(R_i)$ from Eq.~\eqref{eq:range}; weights from
Eqs.~\eqref{eq:wi},\eqref{eq:atten}; normalise $\sum_i w_i=1$.
\item Write \texttt{data/sobp\_layers.csv} with columns
\texttt{energy\_MeV,\,weight}, and \texttt{data/sobp\_design.png}, a sanity plot
of the weights and the analytic SOBP (idealised peaks $+$ attenuation $+$ range
straggling).
\item The Geant4 gun loads the table (\texttt{/stageA/beam/layers}) and samples
one layer per primary.
\end{enumerate}

\section{Verification}
\label{sec:verify}

We run the proton transport with the SOBP gun and inspect the \emph{simulated}
depth--dose (\texttt{data/depth\_dose.csv}). \texttt{field\_design/plot\_sobp.py}
renders it to \texttt{data/sobp\_g4.png} and reports the plateau flatness;
\texttt{analysis\_transport/bragg\_profile.py} is a terminal quick-look.
Acceptance: the dose is flat across $[\Rp,\Rd]$ with a sharp distal fall-off at
$\Rd$. The standard field achieves a flat plateau ($\sim$\SI{7}{\percent} over
the target) with a sharp distal edge. If it were not flat enough, the recourse is
to retune $\mu$, increase $N$, or fall back to the least-squares weights of
Sec.~\ref{sec:weights}.

\section{Lateral spreading, normalisation, and cross-check}
\label{sec:rest}

The depth field above is completed by two further steps, both implemented:
\begin{itemize}[nosep]
\item \textbf{Lateral spreading.} A uniform fluence over the target
cross-section (diameter $\varnothing_{\mathrm t}$) turns the depth plateau into a
dose box. The gun samples primaries uniformly over a disk
(\texttt{/stageA/beam/disk}) matched to the target radius.
\item \textbf{Dose normalisation.} The proton run scores the dose in a target box
inside the phantom; the production is scaled to a clinical dose $D$ via that
target dose ($N_p(\SI{1}{Gy})$ in \texttt{run\_meta.csv}). This fixes the
absolute $\beta^{+}$ production $P_j$ that the decay kinetics consume
(\texttt{03\_decay\_kinetics.tex}).
\end{itemize}
\textbf{Cross-check.} With the field normalised to \SI{1}{Gy}, the integral
$\beta^{+}$ yields are compared to Parodi Table~2
($\sim$\SI{1.8e8}{} decays per Gy, head) by
\texttt{analysis\_transport/parodi\_cross\_check.py}; our homogeneous-brain
yields run about $2.2\times$ the published head values, the expected level of
agreement for a Geant4 cross-section model.

\begin{thebibliography}{9}
\bibitem{bortfeld1996} T.~Bortfeld and W.~Schlegel,
\emph{An analytical approximation of depth--dose distributions for therapeutic
proton beams}, Phys.\ Med.\ Biol.\ \textbf{41}, 1331 (1996); see also
T.~Bortfeld, Med.\ Phys.\ \textbf{24}, 2024 (1997).
\end{thebibliography}

\end{document}
